\documentclass[11pt]{article}
\usepackage[T1]{fontenc}
\usepackage[spanish]{babel}
\usepackage{amsmath,amssymb}
\usepackage[a4paper,margin=2.3cm]{geometry}
\usepackage{graphicx}
\usepackage{tikz}
\usetikzlibrary{arrows.meta,patterns.meta}

\IfFileExists{campo_y_densidad_anillo_matplotlib.pdf}{%
  \newcommand{\includevisualizacion}[2][]{\includegraphics[##1]{##2.pdf}}%
}{%
  \usepackage[inkscapelatex=false]{svg}
  \newcommand{\includevisualizacion}[2][]{\includesvg[##1]{##2}}%
}

\spanishdecimal{.}
\setlength{\parindent}{0pt}
\setlength{\parskip}{0.35em}
\setlength{\abovedisplayskip}{4pt}
\setlength{\belowdisplayskip}{4pt}
\setlength{\abovedisplayshortskip}{4pt}
\setlength{\belowdisplayshortskip}{4pt}

\newcommand{\encabezado}[3]{%
  {\Large\bfseries #1\par}
  \vspace{2pt}
  {\normalsize #2\hfill #3\par}
  \vspace{6pt}\hrule\vspace{8pt}}

\newcommand{\problema}[2]{%
  \vspace{0.55em}
  \noindent\textbf{#1.}\enspace{\small #2}\par
  \vspace{0.3em}}

\newcommand{\inciso}[1]{%
  \vspace{0.35em}
  \noindent\textbf{(#1)}\enspace}

\begin{document}

\encabezado{Campo y potencial electrostático de un anillo con densidad no uniforme}
{Electrodinámica I}{P.S., J.A., J.R.}

\problema{1}{Considere un anillo delgado de radio $R$, situado en el plano
$xy$ y centrado en el origen. El anillo posee una densidad de carga lineal que
varía con la coordenada angular $\varphi'$ de acuerdo con
\begin{equation}
\lambda(\varphi')=\lambda_0\cos\varphi',
\label{eq:densidad}
\end{equation}
donde $\lambda_0$ es una constante positiva con unidades de carga por unidad de
longitud. Se pide:
\begin{enumerate}
  \item \textbf{Discusión sobre la ley de Gauss:} razone de manera argumentada
  por qué esta ley no permite calcular directamente el campo eléctrico de la
  distribución.
  \item \textbf{Integración directa (campo eléctrico):} utilice la ley de
  Coulomb en forma integral para determinar $\vec E$ en un punto arbitrario del
  eje geométrico del anillo, es decir, del eje $z$.
  \item \textbf{Potencial eléctrico y verificación:} calcule $\phi$ sobre el
  eje $z$, tomando $\phi(\infty)=0$, y compruebe la consistencia de la componente
  longitudinal mediante $E_z=-d\phi/dz$.
\end{enumerate}
}

\begin{figure}[htbp]
\centering
\begin{tikzpicture}[
    scale=0.90,
    transform shape,
    line cap=round,
    line join=round,
    eje/.style={-{Stealth[length=1.9mm]},semithick},
    radio/.style={{Stealth[length=1.4mm]}-{Stealth[length=1.4mm]}},
    every node/.style={font=\small}
]
\begin{scope}[x={(1cm,0cm)},y={(0.43cm,0.23cm)},z={(0cm,1cm)}]
\def\Rext{2.22}
\def\Rint{1.78}

% Mitad con densidad positiva: x>0.
\path[
  pattern={Lines[angle=45,distance=3.0pt,line width=0.23pt]},
  pattern color=black!58
]
  plot[domain=-90:90,samples=70,variable=\t]
       ({\Rext*cos(\t)},{\Rext*sin(\t)},0)
  --plot[domain=90:-90,samples=70,variable=\t]
       ({\Rint*cos(\t)},{\Rint*sin(\t)},0)--cycle;

% Mitad con densidad negativa: x<0.
\path[
  pattern={Lines[angle=135,distance=3.0pt,line width=0.23pt]},
  pattern color=black!58
]
  plot[domain=90:270,samples=70,variable=\t]
       ({\Rext*cos(\t)},{\Rext*sin(\t)},0)
  --plot[domain=270:90,samples=70,variable=\t]
       ({\Rint*cos(\t)},{\Rint*sin(\t)},0)--cycle;

% Contornos del anillo.
\draw[very thick]
  plot[domain=0:360,samples=120,variable=\t]
       ({\Rext*cos(\t)},{\Rext*sin(\t)},0);
\draw[very thick]
  plot[domain=0:360,samples=120,variable=\t]
       ({\Rint*cos(\t)},{\Rint*sin(\t)},0);

% Sistema de referencia.
\coordinate (Oe) at (0,0,0);
\draw[eje] (-3.05,0,0)--(3.15,0,0) node[below] {$x$};
\draw[eje] (0,-3.00,0)--(0,3.10,0) node[above right=-1pt] {$y$};
\draw[eje] (0,0,-0.35)--(0,0,3.25) node[left] {$z$};
\fill (Oe) circle (1.1pt);
\node[below left] at (Oe) {$O$};

% Radio fijo. Los dos tramados distinguen los signos de la densidad.
\draw[radio] (Oe)--({2.00*cos(-38)},{2.00*sin(-38)},0)
  node[pos=.58,below=2pt,fill=white,inner sep=1pt] {$R$};

\node[anchor=west,align=left,fill=white,inner sep=2pt]
  at (3.45,0,1.55)
  {$\lambda(\varphi')=\lambda_0\cos\varphi'$\\[1pt]
   $\lambda>0$ para $x>0$\\[-1pt]
   $\lambda<0$ para $x<0$};
\end{scope}
\end{tikzpicture}

\caption{Anillo de radio $R$ situado en el plano $xy$. La proyección oblicua
permite visualizar el eje $z$, mientras los dos tramados distinguen las regiones
con $\lambda>0$ y $\lambda<0$.}
\label{fig:geometria-anillo}
\end{figure}

\textbf{Solución:}

\textbf{1. Discusión sobre la no aplicabilidad de la ley de Gauss.}\par
La ley de Gauss en forma integral establece
\begin{equation}
\oint_{S}\vec E\cdot d\vec S=\frac{q_V}{\varepsilon_0},
\qquad S=\partial V.
\label{eq:gauss}
\end{equation}
Esta relación es válida para cualquier superficie cerrada. Sin embargo, para
despejar el campo a partir del flujo se necesita suficiente simetría: la dirección
del campo debe conocerse de antemano y su magnitud debe ser constante sobre las
porciones relevantes de una superficie gaussiana.

La geometría del anillo es invariante bajo rotaciones alrededor del eje $z$, pero
la distribución de carga no lo es, pues $\lambda$ depende de $\varphi'$. En
particular, la mitad situada en $x>0$ tiene carga positiva y la mitad situada en
$x<0$ tiene carga negativa. Equivalentemente, $\lambda>0$ para
$-\pi/2<\varphi'<\pi/2$, $\lambda<0$ para
$\pi/2<\varphi'<3\pi/2$, y $\lambda=0$ en los extremos que separan ambos
semianillos. Por ello, una esfera o un cilindro coaxial no presenta
un campo de magnitud constante y normal a su superficie. La ley de Gauss permite
conocer el flujo total, pero no simplifica el cálculo directo de $\vec E$.

Además, la carga total del anillo es nula:
\[
q_C=\int_C\lambda\,d\ell'
=\lambda_0R\int_0^{2\pi}\cos\varphi'\,d\varphi'=0.
\]
Esto no significa que el campo sea nulo en todos los puntos; solo indica que el
flujo a través de una superficie que encierre el anillo completo es cero.

\textbf{2. Cálculo del campo eléctrico por integración directa.}\par
Para calcular el campo mediante integración directa partimos de la ley de
Coulomb para una distribución lineal:
\begin{equation}
\vec E(\vec x)=\frac{1}{4\pi\varepsilon_0}
\int_C\lambda(\vec x')
\frac{\vec x-\vec x'}{\lvert\vec x-\vec x'\rvert^3}\,d\ell'.
\label{eq:coulomb}
\end{equation}

\paragraph{Elección de coordenadas y variables.}
Reservamos $\phi$ para el potencial eléctrico y empleamos $\varphi$ como
coordenada angular. El punto de observación $P$ permanece fijo durante la
integración, mientras que el punto fuente $Q$ recorre el anillo. Con el origen en
el centro del anillo y el eje $z$ perpendicular a su plano, escribimos
\begin{equation}
\begin{aligned}
P:&\quad \vec x=z\,\hat z
&&\Longleftrightarrow&&(0,0,z),\\
Q:&\quad \vec x'=R\cos\varphi'\,\hat x
             +R\sin\varphi'\,\hat y
&&\Longleftrightarrow&&(R,\varphi',0).
\end{aligned}
\label{eq:puntos}
\end{equation}
Aquí $z$ fija el punto donde se desea conocer el campo y actúa como parámetro; no
se integra respecto de él. En cambio, $\varphi'$ es la variable de integración y
mapea la distribución de carga. Puesto que el anillo tiene radio fijo $R$,
\[
d\ell'=R\,d\varphi',
\qquad 0\leq\varphi'<2\pi.
\]

La figura~\ref{fig:integracion-anillo} muestra esta elección y los vectores que
intervienen en la ley de Coulomb.

\begin{figure}[htbp]
\centering
\begin{tikzpicture}[
    scale=0.94,
    transform shape,
    line cap=round,
    line join=round,
    eje/.style={-{Stealth[length=2.0mm]},semithick},
    vector/.style={-{Stealth[length=2.4mm]},very thick},
    auxiliar/.style={densely dashed,black!58},
    every node/.style={font=\small}
]
\begin{scope}[x={(1cm,0cm)},y={(0.44cm,0.23cm)},z={(0cm,1cm)}]
\def\Rext{2.28}
\def\Rint{1.84}
\def\angQ{34}
\def\zP{3.65}

\coordinate (O) at (0,0,0);
\coordinate (Q) at ({2.06*cos(\angQ)},{2.06*sin(\angQ)},0);
\coordinate (P) at (0,0,\zP);

% La distribución conserva el mismo código de tramados de la figura 1.
\path[
  pattern={Lines[angle=45,distance=3.1pt,line width=0.23pt]},
  pattern color=black!55
]
  plot[domain=-90:90,samples=70,variable=\t]
       ({\Rext*cos(\t)},{\Rext*sin(\t)},0)
  --plot[domain=90:-90,samples=70,variable=\t]
       ({\Rint*cos(\t)},{\Rint*sin(\t)},0)--cycle;
\path[
  pattern={Lines[angle=135,distance=3.1pt,line width=0.23pt]},
  pattern color=black!55
]
  plot[domain=90:270,samples=70,variable=\t]
       ({\Rext*cos(\t)},{\Rext*sin(\t)},0)
  --plot[domain=270:90,samples=70,variable=\t]
       ({\Rint*cos(\t)},{\Rint*sin(\t)},0)--cycle;
\draw[very thick]
  plot[domain=0:360,samples=120,variable=\t]
       ({\Rext*cos(\t)},{\Rext*sin(\t)},0);
\draw[very thick]
  plot[domain=0:360,samples=120,variable=\t]
       ({\Rint*cos(\t)},{\Rint*sin(\t)},0);

% Sistema de referencia y puntos.
\draw[eje] (O)--(3.70,0,0) node[below] {$x$};
\draw[eje] (O)--(0,3.35,0) node[above right=-1pt] {$y$};
\draw[eje] (O)--(0,0,4.42) node[left] {$z$};
\fill (O) circle (1.1pt);
\node[below left] at (O) {$O$};
\fill (Q) circle (2.2pt);
\node[above right=2pt,fill=white,inner sep=1pt] at (Q) {$Q$};
\draw[thick,fill=white] (P) circle (2pt);
\node[above right=-1pt] at (P) {$P$};

% Vectores de posición y separación.
\draw[vector] (O)--(P)
  node[pos=.62,left=3pt,fill=white,inner sep=1pt]
  {$\vec x=z\,\hat z$};
\draw[vector,densely dashed] (O)--(Q)
  node[pos=.52,above=4pt,fill=white,inner sep=1pt] {$\vec x'$};
\draw[vector,line width=1.1pt] (Q)--(P)
  node[pos=.52,right=3pt,fill=white,inner sep=1pt]
  {$\vec x-\vec x'$};

% Ángulo fuente y elemento de línea curvo.
\draw[-{Stealth[length=1.5mm]}]
  plot[domain=0:\angQ,samples=24,variable=\t]
       ({0.83*cos(\t)},{0.83*sin(\t)},0);
\node[below=4pt,fill=white,inner sep=1pt]
  at ({0.90*cos(17)},{0.90*sin(17)},0) {$\varphi'$};

\draw[line width=2.7pt]
  plot[domain=29:39,samples=18,variable=\t]
       ({2.29*cos(\t)},{2.29*sin(\t)},0);
\draw[-{Stealth[length=1.6mm]},semithick]
  ({2.65*cos(39)},{2.65*sin(39)},0)
  --({2.34*cos(37)},{2.34*sin(37)},0);
\node[anchor=west,align=left,fill=white,inner sep=2pt]
  at ({3.00*cos(42)},{3.00*sin(42)},0)
  {$d\ell'=R\,d\varphi'$\\[-1pt]
   $\lambda(\varphi')=\lambda_0\cos\varphi'$};

\node[anchor=east,align=right,fill=white,inner sep=2pt]
  at (-3.15,0,2.50)
  {$\vec x'$ recorre el anillo;\\
   $\vec x$ permanece fijo};
\end{scope}
\end{tikzpicture}
\caption{Esquema tridimensional para la integración directa. Los tramados
distinguen el signo de $\lambda$. El punto fuente $Q$ recorre el anillo mediante
$\varphi'$, mientras $P$ permanece fijo sobre el eje $z$; el segmento grueso es
el elemento de línea curvo $d\ell'$.}
\label{fig:integracion-anillo}
\end{figure}

De la elección anterior se obtiene
\begin{equation}
\begin{aligned}
\vec x-\vec x'
&=-R\cos\varphi'\,\hat x-R\sin\varphi'\,\hat y+z\,\hat z,\\
\lvert\vec x-\vec x'\rvert
&=\sqrt{R^2\cos^2\varphi'+R^2\sin^2\varphi'+z^2}
=\sqrt{R^2+z^2}.
\end{aligned}
\label{eq:separacion}
\end{equation}
La distancia entre $P$ y cualquier punto del anillo es constante, aunque la
densidad de carga no lo sea. Sustituyendo la ecuación~\eqref{eq:densidad},
$d\ell'=R\,d\varphi'$ y la ecuación~\eqref{eq:separacion} en
\eqref{eq:coulomb}, resulta
\begin{equation}
\vec E(z)=\frac{\lambda_0R}{4\pi\varepsilon_0(R^2+z^2)^{3/2}}
\int_0^{2\pi}\cos\varphi'
\left(-R\cos\varphi'\,\hat x
-R\sin\varphi'\,\hat y+z\,\hat z\right)d\varphi'.
\label{eq:campo-integral}
\end{equation}

Separamos ahora la integral en sus tres componentes. Las integrales angulares
necesarias son
\[
\int_0^{2\pi}\cos^2\varphi'\,d\varphi'=\pi,
\qquad
\int_0^{2\pi}\cos\varphi'\sin\varphi'\,d\varphi'=0,
\qquad
\int_0^{2\pi}\cos\varphi'\,d\varphi'=0.
\]
Por tanto,
\begin{equation}
\begin{aligned}
E_x(z)
&=-\frac{\lambda_0R^2}
{4\pi\varepsilon_0(R^2+z^2)^{3/2}}
\int_0^{2\pi}\cos^2\varphi'\,d\varphi'
=-\frac{\lambda_0R^2}
{4\varepsilon_0(R^2+z^2)^{3/2}},\\[3pt]
E_y(z)
&=-\frac{\lambda_0R^2}
{4\pi\varepsilon_0(R^2+z^2)^{3/2}}
\int_0^{2\pi}\cos\varphi'\sin\varphi'\,d\varphi'=0,\\[3pt]
E_z(z)
&=\frac{\lambda_0Rz}
{4\pi\varepsilon_0(R^2+z^2)^{3/2}}
\int_0^{2\pi}\cos\varphi'\,d\varphi'=0.
\end{aligned}
\label{eq:componentes}
\end{equation}
El campo sobre el eje geométrico del anillo queda entonces
\begin{equation}
\boxed{
\vec E(0,0,z)
=-\frac{\lambda_0R^2}
{4\varepsilon_0(R^2+z^2)^{3/2}}\,\hat x.}
\label{eq:campo-final}
\end{equation}
La dirección $-\hat x$ es consistente con la distribución: el semianillo
positivo repele una carga de prueba hacia $x<0$, mientras que el semianillo
negativo la atrae en la misma dirección. Las componentes en $\hat y$ y $\hat z$
se anulan al integrar el anillo completo. En este sentido, la distribución tiene
carácter dipolar y está orientada a lo largo del eje $x$: aunque el punto de
observación se encuentre sobre el eje $z$, el campo resultante es puramente
transversal y apunta en la dirección $-\hat x$.

\textbf{3. Cálculo del potencial eléctrico y verificación.}\par
El potencial eléctrico de una distribución lineal de carga es
\begin{equation}
\phi(\vec x)=\frac{1}{4\pi\varepsilon_0}
\int_C\frac{\lambda(\vec x')}
{\lvert\vec x-\vec x'\rvert}\,d\ell',
\qquad \phi(\infty)=0.
\label{eq:potencial}
\end{equation}
Sobre el eje $z$, la distancia $\sqrt{R^2+z^2}$ puede extraerse de la integral:
\begin{equation}
\begin{aligned}
\phi(0,0,z)
&=\frac{\lambda_0R}
{4\pi\varepsilon_0\sqrt{R^2+z^2}}
\int_0^{2\pi}\cos\varphi'\,d\varphi'\\
&=0.
\end{aligned}
\label{eq:potencial-eje}
\end{equation}
El potencial es nulo en todos los puntos del eje porque, para cada elemento con
densidad positiva, existe un elemento diametralmente opuesto a la misma distancia
con densidad negativa.

Finalmente, recordamos la relación
\begin{equation}
\vec E=-\nabla\phi.
\label{eq:gradiente}
\end{equation}
Al restringir el potencial al eje se comprueba la componente longitudinal:
\[
E_z(0,0,z)=-\frac{d}{dz}\phi(0,0,z)=0,
\]
en acuerdo con la ecuación~\eqref{eq:componentes}. No hay contradicción entre
$\phi(0,0,z)=0$ y el campo transversal no nulo de
\eqref{eq:campo-final}: la componente $E_x$ depende de cómo varía el potencial al
apartarse del eje,
\[
E_x(0,0,z)=-\left.\frac{\partial\phi}{\partial x}\right|_{(0,0,z)},
\]
información que no está contenida en la función restringida
$\phi(0,0,z)$.

\clearpage
\section*{Visualización numérica}

La figura~\ref{fig:campo-matplotlib} muestra, para $R=1$, la densidad
$\lambda/\lambda_0$ y el campo adimensional
$\widetilde{\vec E}=(4\pi\varepsilon_0R/\lambda_0)\vec E$. Se emplean dos
escalas cromáticas independientes: una divergente para distinguir el signo y
la magnitud de la densidad, y otra logarítmica para la intensidad del campo.
El tamaño de los puntos sobre el anillo se mantiene constante, de modo que toda
la información sobre $\lambda$ queda contenida en el color.

\begin{figure}[htbp]
\centering
\includevisualizacion[width=0.96\textwidth]{campo_y_densidad_anillo_matplotlib}
\caption{Visualización numérica del anillo. En (a), el color representa
$\lambda(\varphi')/\lambda_0$. En (b), las líneas muestran la dirección del
campo en $y=0$ y su color indica $\lvert\widetilde{\vec E}\rvert$. En (c), la
misma información se representa en tres dimensiones: el anillo conserva la
escala de densidad y las flechas utilizan la escala de intensidad del campo.}
\label{fig:campo-matplotlib}
\end{figure}

\end{document}
